% Source: https://artofproblemsolving.com/community/c3h2309983_mathcounts_prep
% Stable identity: aops-topic-2309983-q01
Imagine testing $10000$ people.
\begin{itemize}
\item $2\%$ of them, $200$ people, have the disease, and $95\%$ of these, $190$, test positive.
\item The other $9800$ do not, and $5\%$ of these, $490$, also test positive.
\end{itemize}
So $190+490=680$ people test positive, of whom $190$ have the disease. The probability is
\[\dfrac{190}{680}=\boxed{\dfrac{19}{68}}\approx27.9\%.\]
(By Bayes' theorem: $\dfrac{0.02\cdot0.95}{0.02\cdot0.95+0.98\cdot0.05}=\dfrac{0.019}{0.068}$.)

\textbf{Why so low.} The disease is rare, so the healthy group is $49$ times as large as the sick group. Even a small error rate ($5\%$) applied to this large group produces more false positives ($490$) than there are true positives ($190$), so most positive results come from healthy people.
